\documentclass[12pt]{article}
\usepackage[a4paper, margin=1in]{geometry}
\usepackage{hyperref}

\begin{document}

\begin{center}
    \textbf{\Large PROPOSAL OF EXAMINATION PROJECT} \\
    \vspace{0.2cm}
    \textbf{\large AI Models for Physics 2025/2026}
\end{center}

\vspace{0.5cm}

\noindent
\textbf{Name:} Pallab \hspace{1cm} \textbf{Last name:} Mondal \hspace{1cm} \textbf{Student ID:} 946804 \hspace{1cm} \textbf{Date:} 16/04/2026

\vspace{0.5cm}

\noindent
\textbf{Title:} Physics-Informed ML and Domain Adaptation for Gravitational Lensing

\vspace{0.5cm}

\noindent
\textbf{Learning method/algorithm used:} \\
Physics-Informed Neural Networks (LensPINN), Vision Transformers (HEALSwin), Unsupervised Domain Adaptation (ADDA).

\vspace{0.5cm}

\noindent
\textbf{Objective (max 3 lines):} \\
To build a physics-informed model classifying 4-class dark matter substructures in gravitational lenses, adapting from synthetic dataset simulations to real telescope imagery (HSC) using unsupervised domain adaptation.

\vspace{0.5cm}

\noindent
\textbf{Based on a previously existing project:} \\
( \ ) No, it is totally new not based on code of others \\
(X) Yes, I have partly used code from the following repositories: \\
url: \url{https://github.com/ML4SCI/DeepLense}

\vspace{0.5cm}

\noindent
\textbf{Reference used:} \\
\textbf{Link(s):} \\
\url{http://ml4physicalsciences.github.io/2024/files/NeurIPS_ML4PS_2024_78.pdf} (LensPINN) \\
\url{http://ml4physicalsciences.github.io/2025/files/NeurIPS_ML4PS_2025_252.pdf} (HEAL-PINN) \\
\url{http://arxiv.org/abs/2112.12121} (Domain Adaptation) \\
\url{http://arxiv.org/abs/1702.05464} (ADDA)

\vspace{0.5cm}

\noindent
\textbf{Needed dataset:} \\
( \ ) Not needed \\
( \ ) Yes: I have original data from a research project \\
(X) Yes: Published here url: \url{https://github.com/ML4SCI/DeepLense} (DeepLense synthetic datasets provided by ML4SCI)

\end{document}
